\chapter*{Preface}
\addcontentsline{toc}{chapter}{Preface}
\markboth{Preface}{Preface}

Every phone, car, medical instrument and industrial controller you will work with as an electrical or computer engineer contains
digital logic. Digital design is therefore not a specialist topic: it is the common language in which hardware is described,
verified and built. This book teaches that language from the ground up, starting with a single switch and finishing with a
small general-purpose processor that you will design, simulate and run yourself.

\subsection*{Why this book}
This text grew out of the lectures I give for the courses \emph{\CourseName} (\CourseCode) and \emph{\CourseNameB} (\CourseCodeB) at the
\AuthorDept, \AuthorUni. For several semesters my students asked for one document that follows the lectures in order,
explains the reasoning that slides cannot carry, and offers enough worked examples and exercises to practise on. This book is that document.
It does not replace the classic references listed in the Bibliography; it condenses them into the path that I have found works
for a single semester, and it adds the thing I consider most important in a modern course: every technique is shown
both on paper \emph{and} in VHDL, so that you see the same circuit as a truth table, a schematic, a state diagram and a hardware description.

\subsection*{What is in the book}
The ten chapters follow the course syllabus.
\begin{itemize}
  \item \textbf{Chapters~1--3} introduce digital systems, Boolean algebra, logic gates, logic synthesis and Karnaugh-map optimisation.
  \item \textbf{Chapters~4--5} cover number representation, arithmetic circuits and the standard combinational building blocks (decoders, multiplexers, comparators, ALU).
  \item \textbf{Chapters~6--8} move to sequential logic: latches and flip-flops, registers and counters, and finite state machines (FSMs) in the Moore and Mealy styles.
  \item \textbf{Chapters~9--10} combine everything into processors: first \emph{dedicated} microprocessors (a datapath controlled by an FSM), then a small \emph{general-purpose} processor with an instruction set, an instruction cycle and a program in memory.
\end{itemize}
Four appendices provide a VHDL quick reference, number-system and Boolean-algebra reference tables, the complete listings of the processor designs,
and answers to selected exercises. A glossary and an index are included at the end.

\subsection*{How the book is organised}
Each chapter opens with learning objectives and closes with a summary, exercises and suggestions for further reading.
Worked examples are boxed and numbered; definitions and key design rules are set apart so that they are easy to find again when you revise.
Figures are numbered by chapter, and the VHDL listings can be downloaded together with the book source so that you can run them.

\subsection*{A note on originality and sources}
Digital design rests on well-established results by many people. I have tried to give credit wherever a result, a
historical fact or a data table originates elsewhere, and the Bibliography lists the books and papers I relied on. The explanations,
examples, diagrams, tables and VHDL designs are my own, written for this course.

\subsection*{Feedback}
No book of this size is free of errors. If you find one, or a passage that confused you, please tell me at
\href{mailto:ruslee.s@eng.kmutnb.ac.th}{ruslee.s@eng.kmutnb.ac.th}; corrections will be collected for the next edition.

\bigskip
\begin{flushright}
\AuthorFull\\
\AuthorDept\\
\AuthorUni\\
Bangkok
\end{flushright}
\clearpage
